\section{Object definition in the analysis}
\label{sec:ch6_object_definition}

Two levels of object selection are used in the analysis.
Baseline objects are used for object counting and lepton vetoes.
Signal objects satisfy additional identification, isolation or kinematic requirements and are used to construct the observables defining the signal, control and validation regions.
The requirements below are applied to calibrated objects, and candidates rejected by the overlap removal in Section~\ref{sec:ch6_overlap} are excluded from the event selection.

\subsection{Baseline objects}

The baseline definitions for electrons, muons, hadronic tau candidates and jets are summarised in Table~\ref{tab:ch6_baseline_objects}.
The working points are introduced in Sections~\ref{sec:ch6_electrons}--\ref{sec:ch6_btag}.

The impact-parameter quantities are defined in Section~\ref{sec:ch6_tracks}.
The low-$\pT$ muon rejection used for tau candidates is separate from overlap removal with the baseline muon collection.

\medskip
\noindent\begin{minipage}{\textwidth}
    \begin{spacing}{1}
    \centering
    \small
    \captionof{table}[Baseline object definitions]{Baseline object definitions used for event counting and vetoes in the analysis.}
    \label{tab:ch6_baseline_objects}
    \begin{tabular}{@{}p{0.15\textwidth}p{0.80\textwidth}@{}}
        \toprule
        Object & Requirements \\
        \midrule
        Electron & $\pT>10~\GeV$ and $|\eta|<2.47$, excluding $1.37<|\eta|<1.52$.\newline
        \texttt{TightLH} identification and \texttt{Loose\_VarRad} isolation.\newline
        $|d_0|/\sigma_{d_0}<5$ and $|\Delta z_0\sin\theta|<0.5~\mm$.\newline
        Calorimeter object-quality requirements, including rejection of clusters in the affected end-cap high-voltage regions in 2016. \\
        \addlinespace[0.6em]
        Muon & $\pT>10~\GeV$ and $|\eta|<2.5$.\newline
        \texttt{Loose} identification and \texttt{Loose\_VarRad} isolation.\newline
        $|d_0|/\sigma_{d_0}<3$ and $|\Delta z_0\sin\theta|<0.5~\mm$. \\
        \addlinespace[0.6em]
        Hadronic tau & $\pT>20~\GeV$ and $|\eta|<2.5$, excluding $1.37<|\eta|<1.52$.\newline
        One or three tau-decay tracks with total charge $\pm1$.\newline
        \texttt{Loose} RNN tau identification and \texttt{Loose} RNN electron rejection.\newline
        Rejection of candidates within $\Delta R<0.2$ of a muon with $\pT>2~\GeV$, excluding calorimeter-tagged muons from this test. \\
        \addlinespace[0.6em]
        Jet & Particle-flow anti-$k_t$ jets with $R=0.4$, $\pT>20~\GeV$ and $|\eta|<2.5$.\newline
        \texttt{LooseBad} cleaning and NNJVT \texttt{FixedEffPt} within its calibrated kinematic range.\newline
        No forward jet-vertex requirement. \\
        \bottomrule
    \end{tabular}
    \end{spacing}
\end{minipage}

\subsection{Signal objects}

Signal objects are selected from the baseline collections with the following additional requirements:
\begin{itemize}
    \item \textbf{Electrons.} The \texttt{HighPtCaloOnly} isolation requirement is added to the baseline selection. The transverse-momentum threshold is $28~\GeV$ in the top control regions and $200~\GeV$ in the single-lepton regions without a tau candidate.
    \item \textbf{Muons.} \texttt{Medium} identification and \texttt{Tight\_VarRad} isolation are required in addition to the baseline selection. The transverse-momentum thresholds are the same as for signal electrons.
    \item \textbf{Hadronic taus.} \texttt{Tight} RNN identification is required, and the \texttt{Loose} electron-rejection requirement is retained. The transverse-momentum threshold is $200~\GeV$ in the signal regions and $100~\GeV$ in the top control and validation regions.
    \item \textbf{Jets.} The baseline jet identification is retained. Jets used to define the resonant categories are required to have $\pT>250~\GeV$, while a jet with $50<\pT<250~\GeV$ is required in the non-resonant categories. The $b$-jet subset is selected with the GN2v01 85\% working point introduced in Section~\ref{sec:ch6_btag}.
\end{itemize}
The choice of jet entering each category, object multiplicities, trigger requirements and the remaining event selections are specified in Chapter~\ref{chapter:event_selection}.
